\documentclass[a4paper]{article}
\usepackage[pdftex]{graphicx}
\usepackage[utf8]{inputenc}
\usepackage{enumerate}
\usepackage{siunitx}
\sisetup{locale=DE} 
\usepackage{icomma}
\usepackage{amssymb}
\usepackage{href-ul}
\hypersetup{
	colorlinks=true,
	linkcolor=blue,
	urlcolor=blue}
\usepackage{geometry}
\geometry{a4paper, top=15mm, left=15mm, right=15mm, bottom=15mm,
	headsep=10mm, footskip=12mm}

\begin{document}

\begin{enumerate}[1.]
\item {\bf Gewicht und Masse}
\begin{enumerate}[a)]
\item Erkläre, warum man in der Physik zwischen Gewicht und Masse unterscheidet. (Definition, Zusammenhang, Unterschied)
\item Du landest mit Deinem Raumschiff auf einem unbekannten Exoplaneten und sollst dessen Ortsfaktor bestimmen. Als Hilfsmittel stehen ein Körper mit der Masse $m=\SI{150}{\gram}$ und eine elastische Feder mit der Federkonstanten $ D=\SI{24,0}{\newton\over\meter}$ . Durch Anhängen der Masse dehnt sich die Feder um 13,0 cm.\\
Berechne die Kraft, die auf die Feder wirkt und den Ortsfaktor des Planeten.
\end{enumerate}

%Aufgabe 2
\item {\bf Hooke'sches Gesetz, Dichte}

Ein Aluminiumwürfel mit der Dichte $ \varrho_{Al}=\SI{2,7}{\gram\over\centi\meter^3}$ hat eine Kantenlänge von
 $b=\SI{4,5}{\centi\meter}$ . Er wird an eine Feder mit der Federhärte $D=\SI{0,9}{\newton\over\centi\meter}$ angehängt, die er auf eine Gesamtlänge von  $l=\SI{36}{\centi\meter}$ dehnt. \\
Berechne die Masse des Würfels und die Länge der Feder im unbelasteten Zustand.


\item  {\bf Hooke'sches Gesetz}
\begin{enumerate}[a)]
\item Du sollst untersuchen, ob eine Schraubenfeder bis zur Höchstbelastung von $F_{max}=\SI{12}{\newton}$ dem Gesetz von Hooke gehorcht. Wie gehst du vor?
\item Eine Feder hat die Federhärte $ D_1=\SI{2,5}{\newton\over\centi\meter}$, eine zweite hat die Härte  $ D_2=\SI{5,0}{\newton\over \centi\meter}$. 
\begin{itemize}
\item Welche Gesamtdehnung ergibt sich, wenn man die beiden Federn hintereinander hängt und sie mit $F=\SI{10}{\newton}$ belastet?
\item Welche Federhärte $D_3$ hat diese Kombination?
\end{itemize}
\end{enumerate}


%Aufgabe 4
{\bf \item Dichte}

Herr Müller will seine Garagenauffahrt mit quadratischen Betonplatten mit einer Seitenlänge von $\SI{30,0}{\centi\meter}$ auslegen. Die Betonplatten haben eine Stärke von $\SI{5,0}{\centi\meter}$. Er benötigt insgesamt 150 Betonplatten.  Wie oft muss er in den Baumarkt fahren, wenn sein Anhänger eine Nutzlast von $\SI{0,6}{\tonne}$ hat ? $[ \varrho_{Beton}=\SI{2,4}{\gram\over \centi\meter^3} ]$

\end{enumerate} 

\fbox{
	\begin{minipage}{0.5\textwidth}
		Zur Lösung bitte \href{https://www.okuyakl.de/phys/p8hukL405/ll405.pdf}{hier klicken} oder den QR-Code scannen.\\
	Weitere Arbeitsblätter gibt es unter 
	
	\href{https://www.okuyakl.de}{www.okuyakl.de}
	\end{minipage}
	\hfill
	\begin{minipage}{0.4\textwidth}
		\includegraphics[width=1.5 cm]{../../viecher/zwe03}
		\includegraphics[width=3 cm]{qr405}
		\includegraphics[width=2 cm]{../../viecher/afanticon1}
		
	\end{minipage}}

\end{document}%Lösung-------------------------------------------
